% Part: computability
% Chapter: computability-theory
% Section: rice-theorem

\documentclass[../../../include/open-logic-section]{subfiles}

\begin{document}

\olfileid[pt]{cmp}{thy}{rce}
\olsection{Teorema de Rice}

Se você refletir a respeito, verá que as particularidades de
$\fn{Tot}$ não desempenham nenhum papel na demonstração de
\olref[tot]{prop:total}. Projetamos $h(x,y)$ para se comportar como a
função constante $j(y) = 0$ exatamente quando $x$ pertence a $K$; mas,
nessas circunstâncias, poderíamos igualmente tê-la feito se comportar
como qualquer outra função parcialmente computável. Essa observação
nos permite enunciar um teorema mais geral, que diz, grosso modo, que
nenhuma propriedade não trivial de funções computáveis é decidível.

Tenha em mente que $\cfind{0}$, $\cfind{1}$, $\cfind{2}$,~\dots é
nossa enumeração padrão das funções parcialmente computáveis.

\begin{thm}[Teorema de Rice]
  Seja $C$ um conjunto qualquer de funções parcialmente computáveis, e
  seja $A = \Setabs{n}{\cfind{n} \in C}$. Se $A$ é computável, então
  ou $C$ é vazio, ou $C$ é o conjunto de todas as funções parcialmente
  computáveis.
\end{thm}

Um {\em conjunto de índices} é um conjunto $A$ com a propriedade de
que, se $n$ e $m$ são índices que ``calculam'' a mesma função, então
ou ambos, $n$ e $m$, pertencem a $A$, ou nenhum dos dois pertence.
Não é difícil ver que o conjunto~$A$ do teorema tem essa propriedade.
Reciprocamente, se $A$~é um conjunto de índices e $C$~é o conjunto das
funções calculadas por esses índices, então
$A = \Setabs{n}{\cfind{n} \in C}$.

\begin{explain}
Com essa terminologia, o teorema de Rice equivale a dizer que nenhum
conjunto de índices não trivial é decidível. Para entender o que o
teorema diz, é útil enfatizar a distinção entre \emph{programas}
(digamos, em sua linguagem de programação preferida) e as funções que
eles calculam. Certamente há perguntas computáveis sobre programas
(índices), que são objetos sintáticos: este programa tem mais de 150
símbolos? Tem mais de 22 linhas? Tem uma instrução ``while''? A cadeia
``olá, mundo'' aparece alguma vez como argumento de uma instrução
``print''? O teorema de Rice diz que nenhuma pergunta não trivial
sobre o \emph{comportamento} do programa é computável. Isso inclui
perguntas como estas: o programa para na entrada $0$? Ele para em
algum momento? Em algum momento produz um número par como saída?
\end{explain}

\begin{proof}[Demonstração do teorema de Rice]
Suponha que $C$ não seja vazio nem o conjunto de todas as funções
parcialmente computáveis, e seja $A$ o conjunto dos índices das funções
em~$C$. Mostraremos que, se $A$ fosse computável, poderíamos resolver
o problema da parada; portanto, $A$~não é computável.

Sem perda de generalidade, podemos supor que a função $f$ que não está
definida em ponto algum não pertence a $C$ (caso contrário, troque $C$
por seu complemento no argumento abaixo). Seja $g$ uma função qualquer
em~$C$. A ideia é que, se pudéssemos decidir~$A$, poderíamos distinguir
os índices que calculam~$f$ dos índices que calculam~$g$; e então
poderíamos usar essa capacidade para resolver o problema da parada.

Eis como. Usando as funções parcialmente computáveis universais,
podemos definir uma função
\[
h(x,y) \simeq
\begin{cases}
\text{indefinida} & \text{se $\cfind{x}(x) \fundefined$} \\
g(y) & \text{caso contrário.}
\end{cases}
\]
Para calcular $h$, primeiro tentamos calcular $\cfind{x}(x)$; se essa
computação parar, passamos a calcular~$g(y)$; e, se {\em essa}
computação parar, devolvemos a saída. Mais formalmente, podemos escrever
\[
h(x,y) \simeq \Proj{2}{0}(g(y),\fn{Un}(x,x)).
\]
em que $\Proj{2}{0}(z_0, z_1) = z_0$ é a função de projeção de
$2$ argumentos que devolve o $0$-ésimo argumento e é computável.

Então $h$ é uma composição de funções parcialmente computáveis, e o
lado direito está definido e é igual a~$g(y)$ precisamente quando
$\fn{Un}(x,x)$ e $g(y)$ estão ambas definidas.

Observe que, para um~$x$ fixo, se $\cfind{x}(x)$ é indefinida, então
$h(x,y)$ é indefinida para todo~$y$; e, se $\cfind{x}(x)$ está
definida, então $h(x,y) \simeq g(y)$. Logo, para qualquer valor fixo
de~$x$, ou $h(x,y)$ se comporta exatamente como $f$, ou se comporta
exatamente como $g$, e decidir se $\cfind{x}(x)$ está definida equivale
a decidir qual desses dois casos ocorre. Mas isso equivale a decidir
se $h_x(y) \simeq h(x,y)$ pertence ou não a~$C$ e, se $A$ fosse
computável, poderíamos fazer exatamente isso.

Mais formalmente, como $h$~é parcialmente computável, ela é igual à
função $\cfind{e}$ para algum índice~$e$. Pelo teorema $s$-$m$-$n$,
existe uma função recursiva primitiva~$s$ tal que, para cada~$x$,
$\cfind{s(e,x)}(y) = h_x(y)$. Agora temos que, para cada $x$, se
$\cfind{x}(x) \fdefined$, então $\cfind{s(e,x)}$ é a mesma função
que~$g$ e, portanto, $s(e,x)$ pertence a $A$. Por outro lado, se
$\cfind{x}(x) \uparrow$, então $\cfind{s(e,x)}$ é a mesma função
que~$f$ e, portanto, $s(e,x)$ não pertence a~$A$. Em outras palavras,
temos que, para todo~$x$, $x \in K$ se, e somente se,
$s(e,x) \in A$. Se $A$ fosse computável, $K$~também seria, o que é uma
contradição. Logo, $A$ não é computável.
\end{proof}

O teorema de Rice é muito poderoso. O corolário imediato a seguir
mostra algumas aplicações.
\begin{cor}
Os conjuntos a seguir são indecidíveis.
\begin{enumerate}
\item $\Setabs{x}{\text{$17$ pertence à imagem de $\cfind{x}$}}$
\item $\Setabs{x}{\text{$\cfind{x}$ é constante}}$
\item $\Setabs{x}{\text{$\cfind{x}$ é total}}$
\item $\Setabs{x}{\text{sempre que $y < y'$, $\cfind{x}(y) \fdefined$, e
    se $\cfind{x}(y') \fdefined$, então $\cfind{x}(y) < \cfind{x}(y')$}}$
\end{enumerate}
\end{cor}

\begin{proof}
  Todos esses são conjuntos de índices não triviais.
\end{proof}

\end{document}


